\section{What the development proves, and the arithmetic it uses}
\label{sec:foundations}

Diophantine geometry studies polynomial equations through their rational or
integral points, coordinate changes, heights, and descent. A first-order
arithmetic can express a substantial part of this subject if these objects are
represented by finite data. This paper develops that part over $\Q$. The central
results are an exact bounded-height point enumerator; explicit height distortion
under a nondegenerate map of $\PP^1$; a complete algorithm for its rational
preperiodic points; parametrization and quantitative point counts for a conic
with a supplied point; and the complete calculation
\[
   \{(x,y)\in\Q^2:y^2=x^3-x\}
       =\{(-1,0),(0,0),(1,0)\}.
\]
The last result follows from a fully specified finite descent. We also develop
congruence certificates, rational-root bounds, denominators on integral
Weierstrass equations, and a constructive finite-generation descent interface.

These are a selected theory, not a claim to reproduce the entire subject.
In particular the paper does not prove the general Mordell--Weil, Roth, Siegel,
or Faltings theorems. It does not assert their unprovability in $\iPRA$ either.
Each positive theorem below has its own finite arithmetic proof. Classical
background and terminology may be compared with Milne \cite{milne}, Sutherland
\cite{sutherland2013,sutherland2025}, and Silverman \cite{silverman}.

\subsection{The precise base theory}

We use first-order intuitionistic primitive recursive arithmetic, with one
sort $\N$ and equality. There is a symbol for each fixed primitive recursive
definition: zero, successor, projections, composition, and primitive recursion.
The projection symbols satisfy $\pi_i^k(x_0,\ldots,x_{k-1})=x_i$;
the composition symbol for $f$ and $g_1,\ldots,g_r$ satisfies
$c(\vec x)=f(g_1(\vec x),\ldots,g_r(\vec x))$.
For $h(0,\vec x)=f(\vec x)$ and
$h(n+1,\vec x)=g(n,h(n,\vec x),\vec x)$, these two equations are axioms.
The nonlogical axioms also include successor injectivity and $Sx\ne0$.
Induction is the rule
\[
 \frac{\Gamma\vdash A(0,\vec u)\qquad
       \Gamma,A(n,\vec u)\vdash A(n+1,\vec u)}
      {\Gamma\vdash A(t,\vec u)},
 \qquad A\text{ quantifier-free},
\]
with the usual freshness condition on $n$. The logical rules are intuitionistic
natural deduction: introduction and elimination for $\wedge,\vee,\to,\forall,
\exists$, equality substitution, and elimination of $\bot$. There is no
excluded-middle axiom for arbitrary formulas and no induction on arbitrary
formulas. Quantifier-free formulas are decidable because their atomic
equalities have primitive recursive characteristic functions.

Integers, rationals, vectors and polynomial lists are definitional encodings
in this single numerical sort. For example encode $z\ge0$ by $2z$ and
$z<0$ by $-2z-1$. A rational is a pair $(a,b)$ with $a\in\Z$, $b\ge1$,
$\gcd(|a|,b)=1$; addition and multiplication reduce by gcd. Finite sequences
can be encoded by a length and an iterated pairing, with primitive recursive
access functions. Invalid codes have fixed default outputs. Every geometric
assertion carries the corresponding validity predicate.

\subsection{Bounded quantifiers and explicit witnesses}

A bounded search is a recursion through $0,\ldots,B$. It retains the first
successful index and a flag, or reports failure after exactly $B+1$ tests.
A bounded conjunction is computed by Boolean multiplication. Thus a formula
whose quantifiers are all explicitly bounded has a primitive recursive
characteristic function. Its induction is an instance of quantifier-free
induction applied to that function.

Here is the usual elementary justification, including the point relevant to
weak arithmetic. A search state stores either no successful index, or an
index $j\le i$ for which the tested equality is true. The invariant that the
stored index is valid is quantifier-free. For each additional parameter $k$,
the invariant that a successful test at $k\le i$ has set the flag is also
quantifier-free. Induction on $i$, followed by universal introduction on $k$,
proves success and failure correctness. Bounded conjunction is handled by
searching for the first failed test. Nested bounded quantifiers are compiled
one layer at a time. This explains the use of finite maxima, finite
pigeonhole, and bounded lists later without strengthening induction.

There is an essential distinction between a fixed primitive recursive
definition and a variable code for an arbitrary primitive recursive program.
We use fixed algorithms with numerical inputs, including variable dimensions,
degrees and coefficient lists. We do \emph{not} assume a primitive recursive
universal evaluator for all primitive recursive programs.

\subsection{Elementary arithmetic infrastructure}

\begin{lemma}[Finite arithmetic operations]\label{lem:arithmetic}
The following operations and their stated elementary properties are available
in $\iPRA$: integer arithmetic and order; division with remainder by a positive
integer; gcd and Bezout coefficients; bounded divisibility tests; integer
square root; finite products; and prime factorization of a positive integer
by trial division. Coprime factors of a square are squares. If $a,b$ are
coprime, so are $a^r,b^s$ for positive $r,s$.
\end{lemma}
\begin{proof}
Division of $n$ by $d>0$ searches $q\le n$ for the largest $dq\le n$ and sets
$r=n-dq$. Maximality gives $r<d$. The equations $n=dq+r$ and $0\le r<d$
give uniqueness by subtraction. The square-root algorithm searches
$s\le n$ for the largest $s^2\le n$; it gives
$s^2\le n<(s+1)^2$.

Euclid's remainder algorithm can be run for $a+b+1$ stages, keeping the state
unchanged after a zero remainder. Every active stage strictly reduces the
positive second entry, so it stops in this bound. Maintaining two integer
linear-combination rows proves that the terminal positive entry $g$ is an
integer combination of $a,b$. The identity
$\gcd(a,b)=\gcd(b,a-qb)$ follows directly from the two divisibility
implications. It proves that $g$ divides both inputs and every common divisor
divides $g$. Iterating this construction gives Bezout coefficients for a
nonzero finite vector. The case of zero inputs is handled separately.

For factorization, visit the integers $p=2,\ldots,n$, repeatedly divide the
current residual by $p$ for at most $n$ stages, and record the exponent.
The residual-product invariant is an equation. After processing $p$, no
prime at most $p$ divides the residual: a nontrivial repeated division
strictly decreases it, so $n$ stages suffice to remove that divisor.
This invariant is a bounded test. An active divisor is prime:
if it had a smaller prime divisor, that divisor would already have been
removed. Existence of a prime divisor of a number $>1$ is obtained by its
least divisor in the bounded interval $[2,n]$; a factorization of that divisor
would contradict minimality. Euclid's prime-divisor lemma follows from
Bezout: if $p\nmid a$ then $up+va=1$, and multiplication by $b$ shows
$p\mid ab\Rightarrow p\mid b$. This also proves uniqueness of the exponent
of each prime in a product. If coprime $a,b>0$ have square product, each prime
occurs in just one of them, with even exponent; multiplying the corresponding
half powers constructs their square roots. The same argument proves
coprimality of powers. All divisions, trial indices and exponents are bounded
by the input integer or its explicitly given product. All loop invariants
just used are equations or bounded divisibility assertions, hence Boolean
quantifier-free invariants as above.
\end{proof}

\begin{lemma}[Primitive representatives]\label{lem:normalize}
Every nonzero integral vector $v=(v_0,\ldots,v_n)$ has a primitive
representative $\operatorname{norm}(v)$, obtained by dividing by
$g=\gcd(|v_0|,\ldots,|v_n|)$ and choosing the sign so the first nonzero
coordinate is positive. Two such representatives of the same rational
projective point are equal. Normalization is primitive recursive and never
increases the maximum absolute coordinate.
\end{lemma}
\begin{proof}
The construction and the bound follow from $g\ge1$. Suppose primitive
vectors $v,w$ satisfy $v_iw_j=v_jw_i$ for all $i,j$. Choose an index $j$
with $v_j\ne0$, and Bezout coefficients $c_i$ with $\sum_i c_iv_i=1$.
Set $a=\sum_i c_iw_i$. The proportionality equations give $v_ja=w_j$ and
then $v_jw_i=v_ja v_i$, so cancellation gives $w_i=av_i$ for every $i$.
As $w$ is primitive, $|a|=1$. The sign convention forces $a=1$.
Clearing the finitely many denominators of any nonzero rational vector
reduces to the integral case. Choosing $j$ and summing the Bezout equations
are bounded operations.
\end{proof}

We therefore use canonical primitive vectors as the actual codes for
$\PP^n(\Q)$. Quotient types and choices of representatives are unnecessary
in the object theory. A homogeneous polynomial is a finite list of integer
coefficients and exponent vectors of one declared total degree. The zero
polynomial is allowed with any declared degree; subsequent degree bounds
refer to that declared degree. Its value is computed
by bounded products and sums. Homogeneity shows that its zero locus is
unchanged by scaling and normalization.

\subsection{Relation to the inspected ProveIt foundation}

The inspected repository revision is
\texttt{b5a3ad037723c44d69a3d5b28a950c554327a989}.
The files \texttt{PR.lean}, \texttt{BiInterpretation.lean}, and
\texttt{NativeBiInterpretation.lean} were read at that revision, together
with the ConciseIPRA README. The first defines the first-order theory just
described, with function-symbol type \texttt{PR}, term type \texttt{PTm},
and axiom family \texttt{PRAAx}. The other two expose the translation to the
scalar and native presentations. In particular
\texttt{pra\_iff\_native} transfers a derivation of an iPRA formula to a
derivation of its native reading \cite{proveit}.

This is the appropriate foundation to reuse. The new geometry has not yet
been turned into values of its derivation type. The standalone Lean core in
this package verifies selected arithmetic lemmas; Section~\ref{sec:lean}
specifies the additional work needed for deeply embedded derivations. The
repository's archived-draft index was inspected, but no theorem from an
unread archive is used as a proof dependency.
